




%%%%%%%%%%%%%%%%%%%%%%%%%%%%%%%%%%%%%%%%%%%%%%%%%%%%%%%%%%%%%%%%%
%%%%%%%%%%%%%%%%%%%%%%%%%%%%%%%%%%%%%%%%%%%%%%%%%%%%%%%%%%%%%%%%%
%%%%%%%%%%%%%%%%%%%%%%%%%%%%%%%%%%%%%%%%%%%%%%%%%%%%%%%%%%%%%%%%%

\begin{frame}[t]{Running example: Same-Sex Marriage Survey}
Analysis of US state-level support for same-sex marriage\footnote{
    Based on \textcite{lax:2009:gay,kastellec:2010:laxmrp}.}.

\begin{itemize}
    %
    \item \surcol{\textbf{Survey population:}} Five consistently-coded national polls from 2004 
        ($\surcol{\Nsur ={}\LaxNSur}$)
    \item \surcol{\textbf{Response:}}  $\y = 1$ encodes support for same-sex marriage, $\y = 0$ opposition
       or no opinion.
    \item \tarcol{\textbf{Target population:}} US Census 5\% PUMS data for CA from 2000
    %
\end{itemize}

\pause
We'll analyze this data with two methods: \textbf{calibration weighting} (CW) and \textbf{MrP}
(multilevel regression and poststratification).
%
\end{frame}



%%%%%%%%%%%%%%%%%%%%%%%%%%%%%%%%%%%%%%%%%%%%%%%%%%%%%%%
%%%%%%%%%%%%%%%%%%%%%%%%%%%%%%%%%%%%%%%%%%%%%%%%%%%%%%%
%%%%%%%%%%%%%%%%%%%%%%%%%%%%%%%%%%%%%%%%%%%%%%%%%%%%%%%

\begin{frame}[t]{Two methods}

\splitpagenoline{
    \centering
    \surcol{Survey:}
    Observe \surcol{$(\x_i, y_i)$} for \surcol{$i = 1, \ldots, \nsur$}\\
    Let $\surcol{\Ysur = \{\y_1, \ldots, \y_{\Nsur}\}}$.
}{
    \centering
    \tarcol{Target:}
    Observe \tarcol{$\x_j$} for \tarcol{$j = 1, \ldots, \ntar$}.\\
    We want $\tarcol{\mu := \meantar \y_j}$,
    but don't observe $\tarcol{\y_j}$.  
}



\hrulefill

\splitpagereveal{2-}{5-}{
\centering
\textbf{Calibration weighting estimators:}\\[1em]
Choose weights $\surcol{\w_i}$ using $\x$'s, and take\\
$\muhatcw = \surcol{\meansur \w_i y_i}$.
}{
\centering
\textbf{MrP estimators:}\\[1em]
Estimate $\tarcol{\yhat_j}\surcol{(\Ysur)} = \expect{\post}{\m(\theta^\trans \tarcol{\x_j})}$\\
$\muhatmrp = \tarcol{\meantar \yhat_j}\surcol{(\Ysur)}$
}


\splitpagereveal{3-}{6-}{
    \centering
    Dependence %of \surcol{$\muhat_{\cal}$}
    on \surcol{$\y_i$} is clear\\
    % (\surcol{$\w_i$} typically chosen using only $\x$)
}{
    \centering
    Dependence %of \surcol{$\muhat_\mrp$}
    on \surcol{$\y_i$} can be complicated\\
    (E.g.,~via MCMC draws from $\postsur$)
}
%
\vspace{1em}

\splitpagereveal{4-}{7-}{
    \centering
    Weights give interpretable diagnostics:
    %
    \begin{itemize}
        \item Frequentist variability
        \item Regressor balance / External consistency
        \item Subgroup contribution
        %\item Extrapolation
    \end{itemize}
    %
}{
    \centering
    Black box!
    % \\
    % \vspace{1em}
    % $\leftarrow$ Today, we'll open the box and provide
    %     MrP analogues of all these diagnostics
}


\onslide<8->{
\vspace{1em}

When using OLS, $\tarcol{\yhat_j}\surcol{(\Ysur)}$ is linear 
\parencite{gelman:2007:struggles,benmichael:2021:multilevel,chattopadhyay:2023:implied},
so MrP is a CW estimator.

% \footnote{
%     Most existing literature on comparing CW and MrP focus on linear models. For example,
%     \textcite{gelman:2007:struggles,benmichael:2021:multilevel,chattopadhyay:2023:implied}.}


But currently, for (ubiquitous) MrP models like hierarchical 
logistic regression, no such equivalent weights are available.\parencite{gelman:2007:strugglesrejoinder}
}
% \vspace{2em}
% \hrulefill
% \begin{displayquote}
% ``It would also be desirable to use nonlinear methods  ...
% but then it would seem difficult to construct
% even approximately equivalent weights.  Weighting and fully nonlinear models
% would seem to be completely incompatible methods.''
% --- \parencite{gelman:2007:strugglesrejoinder}
% \end{displayquote}
% \vspace{1em}



\end{frame}




%%%%%%%%%%%%%%%%%%%%%%%%%%%%%%%%%%%%%%%%%%%%%%%%%%%%%%%%%%%%%%%%%
%%%%%%%%%%%%%%%%%%%%%%%%%%%%%%%%%%%%%%%%%%%%%%%%%%%%%%%%%%%%%%%%%
%%%%%%%%%%%%%%%%%%%%%%%%%%%%%%%%%%%%%%%%%%%%%%%%%%%%%%%%%%%%%%%%%

\begin{frame}[t]{Running example: Same-Sex Marriage Survey}
Analysis of US state-level support for same-sex marriage\footnote{
    Based on \textcite{lax:2009:gay,kastellec:2010:laxmrp}.}.

\begin{itemize}
    %
    \item \surcol{\textbf{Survey population:}} Five consistently-coded national polls from 2004 
        ($\surcol{\Nsur ={}\LaxNSur}$)
    \item \surcol{\textbf{Response:}}  $\y = 1$ encodes support for same-sex marriage, $\y = 0$ opposition
       or no opinion.
    \item \tarcol{\textbf{Target population:}} US Census 5\% PUMS data for CA from 2000
    %
\end{itemize}

Let's analyze this data with two methods: \textbf{MrP} and \textbf{calibration weighting} (CW).

\pause
MrP computed with \texttt{brms} \parencite{brms}:\\
{
\tiny
$$
\begin{aligned}
\y \sim{}& \texttt{logit}\Big(~\texttt{(1 | race.female) + (1 | age.cat) + 
  (1 | edu.cat) + (1 | age.edu.cat) +}
  \\{}&\qquad\quad\;\; \texttt{(1 | state) + (1 | region) + (1 | poll) +
  p.relig.full + p.kerry.full}~\Big)
\end{aligned}
$$
}
%
CW used raking on the following coarsened regressors:
({\tiny{\texttt{survey::calibrate}}} \textcite{lumley:2025:survey}):
{
\tiny
$$
\texttt{female, edu.cat, age.cat, race.wbh, race.w.female, age.edu.low}
$$
}
%
\pause
%
\begin{minipage}{0.5\textwidth}
\vspace{-1.5em}
$$
\begin{aligned}
    \textrm{Raw survey mean:} && \surcol{\meansur \y_i} ={}& \LaxYbar \\
    \textrm{Raking:} && \muhatcw ={}& \LaxRakingMu \\
    \textrm{MrP:} && \muhatmrp ={}& \LaxMrpMu\\
\end{aligned}
$$
\end{minipage}
\hfill
\begin{minipage}{0.4\textwidth}
\vspace{-1.5em}
    \pause
    \textbf{How are MrP and raking are using the data differently?}
\end{minipage}
%
%
\end{frame}




%%%%%%%%%%%%%%%%%%%%%%%%%%%%%%%%%%%%%%%%%%%%%%%%%%%%%%%
%%%%%%%%%%%%%%%%%%%%%%%%%%%%%%%%%%%%%%%%%%%%%%%%%%%%%%%
%%%%%%%%%%%%%%%%%%%%%%%%%%%%%%%%%%%%%%%%%%%%%%%%%%%%%%%

\begin{frame}[t]{Linear and nonlinear models}

\splitpagenoline{
\centering
Calibration weighting estimators:\\[1em]
$\muhatcw = \surcol{\meansur w_i y_i}$.}{
\centering
MrP estimators:\\[1em]
$\muhatmrp = \tarcol{\meantar \yhat_j}\surcol{(\Ysur)}$
}

\vspace{1em}

\begin{minipage}{0.3\textwidth}
\only<1-4>{
\YpathPlot{}
}
\only<5->{
\YpathRegionPlot{}
}
\end{minipage}
\begin{minipage}{0.3\textwidth}
\onslide<2->{
\CWPathPlot{}
}
\end{minipage}
\begin{minipage}{0.3\textwidth}
\only<3>{
\MrPPathPlot{}
}
\only<4->{
\MrPlewPathPlot{}
}
\end{minipage}


\onslide<5->{
    The slope is a \textit{MrP locally equivalent weight}, or \textit{MrPlew}.
}
\only<6>{
\\[1em]
The weight $\surcol{\w_i^\mrp}$ measures the effect of local perturbations to $\surcol{\y_i}$ on the estimator $\tarcol{\muhat[\mrp]}$.\\[1em]
Our main contribution is to rigorously connect this idea to CW diagnostics!
}
\only<7->{
$$
\begin{aligned}
\surcol{\w_i^\mrp} :={} 
    \nsur \frac{\partial \muhatmrp}{\partial \surcol{\y_i}} 
        \onslide<8->{
        =
        \nsur  \cov{\post}{
            \tarcol{\meantar \m(\theta^\trans \x_j)},
            \surcol{\frac{\partial}{\partial \y_i} \log \p(\y_i \vert \x_i, \theta)}
        }
    }
\end{aligned}
$$
}
%
\only<8>{
Can estimate $\surcol{\w_i^\mrp}$ \textit{without re-running MCMC}
\parencite{gustafson:1996:localposterior,giordano:2018:covariances}.

}

\end{frame}






%%%%%%%%%%%%%%%%%%%%%%%%%%%%%%%%%%%%%%%%%%%%%%%%%%%%%%%
%%%%%%%%%%%%%%%%%%%%%%%%%%%%%%%%%%%%%%%%%%%%%%%%%%%%%%%
%%%%%%%%%%%%%%%%%%%%%%%%%%%%%%%%%%%%%%%%%%%%%%%%%%%%%%%